\newpage
\section*{Зорилго}

Энэхүү дипломын ажлын зорилго нь Java хэлний annotation болон reflection механизмд тулгуурласан суурь түвшний \textbf{Dependency Injection Framework} зохион байгуулж, хэрэгжүүлэх замаар IoC болон DI зарчмын дотоод ажиллагааг судлах, түүний хэрэглээний боломжийг backend хөгжүүлэлтийн орчинд туршин үнэлэхэд оршино.

\section*{Зорилтууд}

Судалгааны зорилгын хүрээнд дараах зорилтуудыг дэвшүүлж байна:
\begin{enumerate}
    \item Framework, IoC, DI, annotation, reflection зэрэг онолын үндсэн ойлголтуудыг судлах;
    \item Spring Framework-ийн dependency injection болон annotation боловсруулах механизмыг шинжлэх;
    \item Java Reflection API ашиглан класс, талбар, гишүүн функцийн мэдээллийг тодорхойлох аргуудыг судлах;
    \item Custom annotation-ууд тодорхойлон dependency-ийг автоматаар илрүүлэх боломжийг бүрдүүлэх;
    \item IoC container-ийн үндсэн бүтэц болон bean удирдлагын механизмыг зохион бүтээх;
    \item суурь түвшний хялбаршуулсан DI framework-ийг хэрэгжүүлж турших;
    \item Хэрэгжүүлсэн framework-ийн ажиллагааг жишээ backend орчинд шалгаж, үр дүнг үнэлэх.
\end{enumerate}

\section*{Судалгааны объект}

Энэхүү судалгааны объект нь Java хэл дээрх annotation болон reflection-д суурилсан dependency удирдлагын механизм бүхий backend framework-ийн дотоод архитектур болно.

\section*{Судлагдахуун}

Судлагдахуун нь IoC container, Dependency Injection механизм, annotation боловсруулах аргачлал, reflection ашиглан объект үүсгэх болон dependency холбох процесс, мөн эдгээрийг framework хэлбэрээр хэрэгжүүлэх зарчим юм.

\section*{Судалгааны арга зүй}
Энэхүү судалгаанд харьцуулсан шинжилгээний аргыг ашиглан Spring Framework-ийн dependency injection болон annotation боловсруулах механизмыг судалж, түүнийг хэрэгжүүлсэн DI framework-ийн бүтэц, ажиллагаатай харьцуулан үнэлнэ. Үүнд объект үүсгэх зарчим, dependency холболт, container-ийн үүрэг, annotation таних болон боловсруулах ажиллагаа, өргөтгөх боломж зэрэг үзүүлэлтүүдийг харгалзан ижил болон ялгаатай талыг тодорхойлно.
Энэхүү судалгаанд дараах үндсэн арга зүйг ашиглана:
\begin{itemize}
    \item \textbf{Баримт бичиг, эх сурвалжийн судалгаа} -- Framework-ийн архитектур, IoC, DI, reflection, annotation-ийн онолын үндсийг судлах;
    \item \textbf{Харьцуулсан шинжилгээ} -- Spring Framework-ийн механизмтай харьцуулан судлах;
    \item \textbf{Загварчлал, зохиомжийн арга} -- DI framework-ийн бүтэц, бүрэлдэхүүн хэсгүүдийг тодорхойлох;
    \item \textbf{Туршилт, хэрэгжүүлэлт} -- Java хэл дээр framework боловсруулж, жишээ систем дээр турших;
    \item \textbf{Үр дүнгийн үнэлгээ} -- хэрэгжүүлсэн framework-ийн ажиллагаа, давуу тал, хязгаарлалтыг үнэлэх.
\end{itemize}

\section*{Судалгааны ажлын шинэлэг тал}

Энэхүү ажлын шинэлэг тал нь өргөн хэрэглэгддэг Spring Framework-ийн дотоод ажиллагааны зарчмыг зөвхөн онолын түвшинд тайлбарлах бус, annotation болон reflection-д суурилсан суурь түвшний \textbf{DI Framework} болгон бие даан хэрэгжүүлж туршсанаар backend framework-ийн үндсэн механизмыг илүү ойлгомжтой, бодит түвшинд харуулахад оршино.

\section*{Судалгааны ажлын практик ач холбогдол}

Энэхүү судалгааны үр дүнд боловсруулсан framework нь IoC болон DI зарчмын хэрэгжилтийг ойлгоход сургалтын болон туршилтын зориулалтаар ашиглагдах боломжтой. Мөн backend хөгжүүлэлтийн суурь архитектур, dependency удирдлага, annotation болон reflection механизмын талаарх мэдлэгийг гүнзгийрүүлэхэд практик ач холбогдолтой. Цаашлаад энэхүү ажлын үр дүн нь framework-ийн дотоод ажиллагаа, шийдлийг судлах,судалгааны дагуу хэрэгжүүлэн сургалтын хэрэглэгдэхүүн болгон ашиглах үндэс суурь болно.

\section*{Дипломын ажлын бүтэц}

Энэхүү дипломын ажил нь үндсэндээ онолын судалгаа, системийн зохиомж, хэрэгжүүлэлт, туршилт, үр дүнгийн үнэлгээ гэсэн хэсгүүдээс бүрдэнэ. Нэгдүгээр бүлэгт судалгааны онолын үндэс, framework болон dependency injection-ийн талаар авч үзнэ. Хоёрдугаар бүлэгт системийн шаардлага шинжилгээг гаргах.Гуравдугаар бүлэгт зохиомж, архитектурын шийдлийг тайлбарлана. Дөрөвдүгээр бүлэгт framework-ийн хэрэгжүүлэлт, туршилт, үр дүнг танилцуулж, эцэст нь дүгнэлт, цаашдын хөгжүүлэлтийн боломжийг тусгана.